\section{Appointed Texts of the Mass}
\runninghead{Appointed Texts of the Mass}

The Latin texts below reproduce the Mass of \latin{Dominica prima Adventus} visually collated against the 1962 typical \work{Missale Romanum} (printed pp.~1--2; Preface of the Most Holy Trinity, n.~1082, printed p.~293) and corroborated word for word in pre-1931 public-domain editions (Pustet Ratisbon 1862, Tournai Descl\'ee 1889, Tours Mame 1922, and Cummiskey 1861), with ligatures expanded and chant accents omitted. The accompanying English texts come from two public-domain historical translations used for study, not as official 1962 liturgical English: the Douay--Rheims Bible (Challoner revision, 1899 American edition, resolved in Vulgate psalm numbering) for biblical texts, and Fr.~F.~X. Lasance and Fr.~Francis Augustine Walsh, O.S.B., \work{The New Roman Missal in Latin and English} (Benziger Brothers, revised 1945) for the three composed orations and the Preface of the Most Holy Trinity. Liturgical adaptations separating the Missal's Latin from the Clementine Vulgate and Douay--Rheims text are noted beside each passage.

\subsection{Introit}\label{proper-introit}
\properlocus{1962 \work{Missale Romanum}, printed p.~1, unit~1 (\latin{Ant.~ad Introitum}): Ps.~24:1--3, 4b (Heb.~25:1--3a, 4). The Missal omits the psalm title (\latin{In finem. Psalmus David.}) and vocative \latin{Domine} in v.~1, retains Old Latin \latin{qui te exspectant} against Gallican \latin{qui sustinent te} in v.~3b, omits the imprecatory half-verse v.~4a (\latin{Confundantur omnes iniqua agentes supervacue.}), and takes v.~4b in Gallican wording (\latin{demonstra mihi}) as the verse. \latin{Gloria Patri} is retained throughout Advent (RGMR~427--429), after which the antiphon is repeated.}

\begin{latinproper}
Ad te levavi animam meam: Deus meus, in te confido, non erubescam: neque irrideant me inimici mei: etenim universi, qui te exspectant, non confundentur. \textnormal{Ps.} Vias tuas, Domine, demonstra mihi: et semitas tuas edoce me. \textnormal{V.} Gloria Patri, et Filio, et Spiritui Sancto. Sicut erat in principio, et nunc, et semper, et in saecula saeculorum. Amen. \textnormal{[Antiphon repeated: Ad te levavi\dots{} non confundentur.]}
\end{latinproper}

\begin{englishwitness}{Douay--Rheims (Challoner, 1899 American edition), Ps.~24:1--3, 4b}
To thee, O Lord, have I lifted up my soul. In thee, O my God, I put my trust; let me not be ashamed. Neither let my enemies laugh at me: for none of them that wait on thee shall be confounded. \textit{V.} Shew, O Lord, thy ways to me, and teach me thy paths.
\end{englishwitness}

\subsection{Collect}\label{proper-collect}
\properlocus{1962 \work{Missale Romanum}, printed p.~1, unit~2 (\latin{Oratio}). Composed Roman prayer attested in the Gregorian Sacramentary (\work{Hadrianum}, p.~113); verbally echoes Ps.~79:3 (\latin{Excita potentiam tuam, et veni}). Addressed directly to Christ the Son (\latin{Qui vivis et regnas\dots}). Sole Collect on Sunday, 29~November~2026 under universal 1962 rubrics. Lasance preserves the abbreviated hand-missal conclusion ``Who livest.'' where the 1962 altar Missal prints the full trinitarian ending.}

\begin{latinproper}
Excita, quaesumus, Domine, potentiam tuam, et veni: ut ab imminentibus peccatorum nostrorum periculis, te mereamur protegente eripi, te liberante salvari: Qui vivis et regnas cum Deo Patre in unitate Spiritus Sancti, Deus: per omnia saecula saeculorum. \textnormal{R.} Amen.
\end{latinproper}

\begin{englishwitness}{Lasance and Walsh, The New Roman Missal (1945), p.~103}
BESTIR, O Lord, Thy might, we pray Thee and come; That, defended by Thee, we may deserve rescue from approaching dangers brought on by our sins, and being set free by Thee, obtain our salvation. Who livest.
\end{englishwitness}

\subsection{Epistle}\label{proper-epistle}
\properlocus{1962 \work{Missale Romanum}, printed pp.~1--2, unit~3 (\latin{Lectio Epistolae beati Pauli Apostoli ad Romanos}): printed citation Rom.~13:11--14; actual liturgical extent Rom.~13:11--14a (57~spoken words). The Missal prepends \latin{Fratres:}, omits Vulgate \latin{Et hoc\dots{} tempus} in v.~11, and closes at \latin{Dominum Jesum Christum.}, omitting v.~14b (\latin{et carnis curam ne feceritis in desideriis}, ``and make not provision for the flesh in its concupiscences''). Public-domain Tours Mame 1922 prints \latin{Abjiciamus}/\latin{commessationibus} where 1962 prints \latin{Abiciamus}/\latin{comessationibus}.}

\begin{latinproper}
Fratres: Scientes quia hora est jam nos de somno surgere. Nunc enim propior est nostra salus, quam cum credidimus. Nox praecessit, dies autem appropinquavit. Abjiciamus ergo opera tenebrarum, et induamur arma lucis. Sicut in die honeste ambulemus: non in commessationibus et ebrietatibus, non in cubilibus et impudicitiis, non in contentione et aemulatione: sed induimini Dominum Jesum Christum.
\end{latinproper}

\begin{englishwitness}{Douay--Rheims (Challoner, 1899 American edition), Rom.~13:11--14a}
And that knowing the season; that it is now the hour for us to rise from sleep. For now our salvation is nearer than when we believed. The night is passed, and the day is at hand. Let us therefore cast off the works of darkness, and put on the armour of light. Let us walk honestly, as in the day: not in rioting and drunkenness, not in chambering and impurities, not in contention and envy: But put ye on the Lord Jesus Christ
\end{englishwitness}

\subsection{Gradual}\label{proper-gradual}
\properlocus{1962 \work{Missale Romanum}, printed p.~2, unit~4 (\latin{Graduale}): Ps.~24:3b, 4b (Heb.~25:3a, 4). The Missal drops \latin{etenim} (``for'') at the start of the respond, retains Old Latin \latin{qui te exspectant}, appends the vocative \latin{Domine}, and sings the Old Latin verb \latin{notas fac mihi} (``make known to me'') in the verse, contrasting with Gallican \latin{demonstra mihi} in the Introit verse of the same Mass.}

\begin{latinproper}
Universi, qui te exspectant, non confundentur, Domine. \textnormal{V.} Vias tuas, Domine, notas fac mihi: et semitas tuas edoce me.
\end{latinproper}

\begin{englishwitness}{Douay--Rheims (Challoner, 1899 American edition), Ps.~24:3b, 4b}
for none of them that wait on thee shall be confounded. \textit{V.} Shew, O Lord, thy ways to me, and teach me thy paths.
\end{englishwitness}

\subsection{Alleluia}\label{proper-alleluia}
\properlocus{1962 \work{Missale Romanum}, printed p.~2, unit~5: Ps.~84:8 (Heb.~85:8; modern English 85:7). Required on Sunday. By the printed 1962 rubric standing immediately before this unit, when the Sunday Mass is resumed on an Advent weekday during the week (RGMR~299), the Alleluia and its verse are omitted and only the Gradual is said, without a Tract.}

\begin{latinproper}
Alleluia, alleluia. \textnormal{V.} Ostende nobis, Domine, misericordiam tuam: et salutare tuum da nobis. Alleluia.
\end{latinproper}

\begin{englishwitness}{Douay--Rheims (Challoner, 1899 American edition), Ps.~84:8 (liturgical Alleluia frame)}
Alleluia, alleluia. \textit{V.} Shew us, O Lord, thy mercy; and grant us thy salvation. Alleluia.
\end{englishwitness}

\subsection{Gospel}\label{proper-gospel}
\properlocus{1962 \work{Missale Romanum}, printed p.~2, unit~6 (\latin{Sequentia sancti Evangelii secundum Lucam}): Luke~21:25--33. The Missal replaces Vulgate opening \latin{Et} with the liturgical opening \latin{In illo tempore: Dixit Iesus discipulis suis:} (``At that time, Jesus said to his disciples:''). At v.~29, the 1899 American Douay--Rheims edition prints ``And he spoke to them in a similitude.'' where the Challoner Gutenberg 1581 edition omits the preposition (``And he spoke to them a similitude.''). Followed on Sundays by the \latin{Credo} (printed rubric, p.~2; RGMR~475--476).}

\begin{latinproper}
In illo tempore: Dixit Iesus discipulis suis: Erunt signa in sole, et luna, et stellis, et in terris pressura gentium prae confusione sonitus maris, et fluctuum: arescentibus hominibus prae timore et exspectatione, quae supervenient universo orbi: nam virtutes caelorum movebuntur. Et tunc videbunt Filium hominis venientem in nube cum potestate magna, et maiestate. His autem fieri incipientibus, respicite, et levate capita vestra: quoniam appropinquat redemptio vestra. Et dixit illis similitudinem: Videte ficulneam, et omnes arbores: cum producunt iam ex se fructum, scitis quoniam prope est aestas. Ita et vos cum videritis haec fieri, scitote quoniam prope est regnum Dei. Amen dico vobis, quia non praeteribit generatio haec, donec omnia fiant. Caelum et terra transibunt: verba autem mea non transibunt.
\end{latinproper}

\begin{englishwitness}{Douay--Rheims (Challoner, 1899 American edition), Luke~21:25--33}
And there shall be signs in the sun, and in the moon, and in the stars; and upon the earth distress of nations, by reason of the confusion of the roaring of the sea and of the waves; Men withering away for fear, and expectation of what shall come upon the whole world. For the powers of heaven shall be moved; And then they shall see the Son of man coming in a cloud, with great power and majesty. But when these things begin to come to pass, look up, and lift up your heads, because your redemption is at hand. And he spoke to them in a similitude. See the fig tree, and all the trees: When they now shoot forth their fruit, you know that summer is nigh; So you also, when you shall see these things come to pass, know that the kingdom of God is at hand. Amen, I say to you, this generation shall not pass away, till all things be fulfilled. Heaven and earth shall pass away, but my words shall not pass away.
\end{englishwitness}

\subsection{Offertory}\label{proper-offertory}
\properlocus{1962 \work{Missale Romanum}, printed p.~2, unit~7 (\latin{Ant.~ad Offertorium}): Ps.~24:1--3 (Heb.~25:1--3a). Identical word for word to the Introit antiphon, retaining Old Latin \latin{exspectant} and omitting the superscription and v.~1 \latin{Domine}.}

\begin{latinproper}
Ad te levavi animam meam: Deus meus, in te confido, non erubescam: neque irrideant me inimici mei: etenim universi, qui te exspectant, non confundentur.
\end{latinproper}

\begin{englishwitness}{Douay--Rheims (Challoner, 1899 American edition), Ps.~24:1--3}
To thee, O Lord, have I lifted up my soul. In thee, O my God, I put my trust; let me not be ashamed. Neither let my enemies laugh at me: for none of them that wait on thee shall be confounded.
\end{englishwitness}

\subsection{Secret}\label{proper-secret}
\properlocus{1962 \work{Missale Romanum}, printed p.~2, unit~8 (\latin{Secreta}). Composed Roman prayer attested identically in the Gregorian Sacramentary (\work{Hadrianum}, p.~113, \latin{Super oblata}). Said silently except the closing \latin{Per omnia saecula saeculorum} (RGMR~480--481). Lasance abbreviates the conclusion to ``Through our Lord.'' where 1962 prints the full trinitarian ending.}

\begin{latinproper}
Haec sacra nos, Domine, potenti virtute mundatos, ad suum faciant puriores venire principium. Per Dominum nostrum Jesum Christum, Filium tuum: Qui tecum vivit et regnat in unitate Spiritus Sancti, Deus: per omnia saecula saeculorum. \textnormal{R.} Amen.
\end{latinproper}

\begin{englishwitness}{Lasance and Walsh, The New Roman Missal (1945), p.~105}
CLEANSING us by their mighty power, may these Holy Mysteries, O Lord, make us come more pure before Thee who art their author. Through our Lord.
\end{englishwitness}

\subsection{Communion}\label{proper-communion}
\properlocus{1962 \work{Missale Romanum}, printed p.~2, unit~9 (\latin{Ant.~ad Communionem}): Ps.~84:13 (Heb.~85:13; modern English 85:12). The Missal omits Vulgate opening \latin{Etenim} (``For''). St.~Augustine's Old Latin Psalter reads \latin{suavitatem} (``sweetness'') where the Gallican Vulgate and Roman Missal sing \latin{benignitatem} (``goodness'').}

\begin{latinproper}
Dominus dabit benignitatem: et terra nostra dabit fructum suum.
\end{latinproper}

\begin{englishwitness}{Douay--Rheims (Challoner, 1899 American edition), Ps.~84:13}
For the Lord will give goodness: and our earth shall yield her fruit.
\end{englishwitness}

\subsection{Postcommunion}\label{proper-postcommunion}
\properlocus{1962 \work{Missale Romanum}, printed p.~2, unit~10 (\latin{Postcommunio}). Composed Roman prayer attested in the Gregorian Sacramentary (\work{Hadrianum}, p.~113, \latin{Ad complendum}, reading \latin{et reparationis} where the 1570--1962 Missals read \latin{ut reparationis}). Verbally echoes Ps.~47:10 (Heb.~48:10; modern English~48:9), \latin{Suscepimus, Deus, misericordiam tuam in medio templi tui}, turning the psalm's perfect tense into the jussive subjunctive \latin{Suscipiamus}. The 1962 Missal prints the cue \latin{Per Dominum nostrum.}; Lasance prints ``Through our Lord.''}

\begin{latinproper}
Suscipiamus, Domine, misericordiam tuam in medio templi tui: ut reparationis nostrae ventura solemnia congruis honoribus praecedamus. Per Dominum nostrum.
\end{latinproper}

\begin{englishwitness}{Lasance and Walsh, The New Roman Missal (1945), p.~106}
MAY we receive Thy mercy, O Lord, in the midst of Thy temple, that with due reverence we may prepare for the coming festival of our redemption. Through our Lord.
\end{englishwitness}

\subsection{Preface of the Most Holy Trinity (Governed Ordinary)}
\properlocus{Appointed for Sundays of Advent by explicit rubric on 1962 \work{Missale Romanum}, printed p.~2 (\latin{Praefatio de sanctissima Trinitate}; text at n.~1082, printed p.~293; RGMR~494~b). On Advent weekdays resuming the Sunday Mass, the Common Preface is said instead (n.~1086, printed p.~297; RGMR~482, 498). Latin text: Cummiskey 1861, collated against 1962 n.~1082.}

\begin{latinproper}
Vere dignum et justum est, aequum et salutare, nos tibi semper, et ubique gratias agere: Domine sancte, Pater omnipotens, aeterne Deus: Qui cum unigenito Filio tuo et Spiritu Sancto unus es Deus, unus es Dominus: non in unius singularitate Personae, sed in unius Trinitate substantiae. Quod enim de tua gloria, revelante te credimus, hoc de Filio tuo, hoc de Spiritu Sancto, sine differentia discretionis sentimus. Ut in confessione verae, sempiternaeque Deitatis, et in Personis proprietas, et in essentia Unitas, et in Majestate adoretur aequalitas. Quam laudant angeli atque archangeli, cherubim quoque ac seraphim: qui non cessant clamare quotidie, una voce dicentes:
\end{latinproper}

\begin{englishwitness}{Lasance and Walsh, The New Roman Missal (1945), pp.~773--774}
IT IS truly meet and just, right and profitable, for us, at all times, and in all places, to give thanks to Thee, O Lord, the holy One, the Father almighty, the everlasting God: Who, together with Thine only-begotten Son and the Holy Ghost, art one God, one Lord, not in the singleness of one Person, but in the Trinity of one substance. For that which, according to Thy revelation, we believe of Thy glory, the same we believe of Thy Son, the same of the Holy Ghost, without difference or distinction; so that in the confession of one true and eternal Godhead we adore distinctness in persons, oneness in essence, and equality in majesty: Which the angels praise, and the archangels, the cherubim also and the seraphim, who cease not, day by day crying out with one voice, to repeat:
\end{englishwitness}
